\chapter{The Category of Fuzzy Domains}
\label{ch:category-fuzzy-domains}

In Chapter~\ref{ch:fuzzy-ideal-monad}, the upper fuzzy character adjunction
produced the fuzzy ideal monad.  After the finite-flat-weight representation and
the strictification convention of that chapter, this monad is written
\[
    T_{\mathbb V}X=\operatorname{Idl}_{\mathbb V}(X),
\]
where \(\operatorname{Idl}_{\mathbb V}(X)\) is the full
\(\mathbb V\)-subcategory of
\[
    [X^{\op},\underline{\mathbb V}]
\]
spanned by the finite-flat weights
\[
    I:X^{\op}\to\underline{\mathbb V}.
\]
Here \(\underline{\mathbb V}\) denotes the self-enriched
\(\mathbb V\)-category of truth values.

The unit is the Yoneda embedding
\[
    \eta_X:X\to\operatorname{Idl}_{\mathbb V}(X),
    \qquad
    x\longmapsto X(-,x),
\]
and the multiplication is flattening:
\[
    \mu_X(\mathcal I)(x)
    =
    \int^{I\in\operatorname{Idl}_{\mathbb V}(X)}
        \mathcal I(I)\tensor I(x).
\]
All occurrences of \(\operatorname{Idl}_{\mathbb V}(X)\), functor categories,
and Eilenberg--Moore categories are interpreted under the universe conventions
fixed earlier.  Thus, when an ideal category is not small in the base universe,
it is formed in the next universe and regarded there as small for the relevant
enriched constructions.

The purpose of this chapter is to study the ordinary Eilenberg--Moore category
of strict algebras for this monad.  These strict algebras are the fuzzy domains.
The main result is that the category of fuzzy domains is symmetric monoidal
closed.  The tensor product is obtained by applying Kelly's tensor-product
construction to the saturated commutative doctrine of finite-flat weighted
colimits, and the internal hom is the fuzzy domain of finite-flat-colimit-
preserving functors.

The point of the chapter is not to add new structure to fuzzy domains.  Rather,
we show that the finite-flat doctrine already used to define fuzzy ideals also
determines the colimits, tensor product, and internal hom of fuzzy domains.

\section{The finite-flat doctrine}
\label{sec:finite-flat-doctrine-fdom}

We first isolate the closure properties of finite-flat weights needed later.

\begin{definition}[Finite-flat weights]
\label{def:finite-flat-weights-fdom}
Let \(X\) be a small \(\mathbb V\)-category.  A weight
\[
    I:X^{\op}\to\underline{\mathbb V}
\]
has associated weighted-colimit \(\mathbb V\)-functor
\[
    \chi_I:[X,\underline{\mathbb V}]\to\underline{\mathbb V},
    \qquad
    \chi_I(U)=\int^{x\in X} I(x)\tensor U(x).
\]
We say that \(I\) is finite-flat if \(\chi_I\) preserves the finite limits of
the fuzzy-frame doctrine.

We write
\[
    \Phi_{\mathrm{ff}}(X)
\]
for the full \(\mathbb V\)-subcategory of
\([X^{\op},\underline{\mathbb V}]\) spanned by the finite-flat weights.  Thus
\[
    \Phi_{\mathrm{ff}}(X)
    =
    \operatorname{Idl}_{\mathbb V}(X).
\]
\end{definition}

\begin{lemma}[Representables are finite-flat]
\label{lem:representables-finite-flat-fdom-rewrite}
For every object \(x\in X\), the representable weight
\[
    X(-,x):X^{\op}\to\underline{\mathbb V}
\]
is finite-flat.
\end{lemma}

\begin{proof}
The colimit functional associated to \(X(-,x)\) is
\[
    \chi_{X(-,x)}(U)
    =
    \int^{y\in X}X(y,x)\tensor U(y).
\]
By the enriched co-Yoneda formula,
\[
    \chi_{X(-,x)}(U)\cong U(x).
\]
Thus \(\chi_{X(-,x)}\) is evaluation at \(x\).  Since the finite limits of the
fuzzy-frame doctrine in \([X,\underline{\mathbb V}]\) are computed pointwise,
evaluation at \(x\) preserves them.  Hence \(X(-,x)\) is finite-flat.
\end{proof}

\begin{lemma}[Finite-flat direct images]
\label{lem:finite-flat-direct-images-fdom-rewrite}
Let
\[
    D:K\to X
\]
be a \(\mathbb V\)-functor, and let
\[
    W:K^{\op}\to\underline{\mathbb V}
\]
be finite-flat.  Define
\[
    D_!W:X^{\op}\to\underline{\mathbb V}
\]
by the lower direct-image formula
\[
    (D_!W)(x)
    =
    \int^{k\in K} W(k)\tensor X(x,Dk).
\]
Then \(D_!W\) is finite-flat.
\end{lemma}

\begin{proof}
Let
\[
    U:X\to\underline{\mathbb V}
\]
be an upper predicate on \(X\).  The colimit functional associated to \(D_!W\)
satisfies
\[
\begin{aligned}
    \chi_{D_!W}(U)
    &=
    \int^{x\in X}(D_!W)(x)\tensor U(x)
    \\
    &=
    \int^{x\in X}
        \left(
            \int^{k\in K}W(k)\tensor X(x,Dk)
        \right)\tensor U(x)
    \\
    &\cong
    \int^{k\in K}
        W(k)\tensor
        \left(
            \int^{x\in X}X(x,Dk)\tensor U(x)
        \right)
    \\
    &\cong
    \int^{k\in K}W(k)\tensor U(Dk)
    \\
    &=
    \chi_W(UD).
\end{aligned}
\]
The first isomorphism is Fubini for coends, and the second is co-Yoneda.

Precomposition
\[
    D^*:[X,\underline{\mathbb V}]\to[K,\underline{\mathbb V}],
    \qquad
    U\longmapsto UD,
\]
preserves the finite limits of the fuzzy-frame doctrine, since these limits are
computed pointwise.  Since \(W\) is finite-flat, \(\chi_W\) preserves those
finite limits.  Hence
\[
    \chi_{D_!W}
    \cong
    \chi_WD^*
\]
preserves them.  Thus \(D_!W\) is finite-flat.
\end{proof}

\begin{proposition}[Finite-flat saturation]
\label{prop:finite-flat-saturation-fdom-rewrite}
Finite-flat weights are closed under finite-flat weighted colimits in presheaf
categories.

Explicitly, let
\[
    W:K^{\op}\to\underline{\mathbb V}
\]
be finite-flat, and let
\[
    H:K\to\operatorname{Idl}_{\mathbb V}(X)
\]
be a \(\mathbb V\)-functor.  Define
\[
    J:X^{\op}\to\underline{\mathbb V}
\]
pointwise by
\[
    J(x)
    =
    \int^{k\in K} W(k)\tensor Hk(x).
\]
Then \(J\) is finite-flat.
\end{proposition}

\begin{proof}
Since presheaf colimits are computed pointwise, \(J\) is the \(W\)-weighted
colimit of the diagram
\[
    H:K\to [X^{\op},\underline{\mathbb V}]
\]
in the ambient presheaf category.  Hence \(J\) is again a weight
\[
    J:X^{\op}\to\underline{\mathbb V}.
\]

For every upper predicate
\[
    U:X\to\underline{\mathbb V},
\]
we have
\[
\begin{aligned}
    \chi_J(U)
    &=
    \int^{x\in X}J(x)\tensor U(x)
    \\
    &=
    \int^{x\in X}
        \left(
            \int^{k\in K}W(k)\tensor Hk(x)
        \right)\tensor U(x)
    \\
    &\cong
    \int^{k\in K}
        W(k)\tensor
        \left(
            \int^{x\in X}Hk(x)\tensor U(x)
        \right)
    \\
    &=
    \int^{k\in K}W(k)\tensor \chi_{Hk}(U).
\end{aligned}
\]
Thus
\[
    \chi_J(U)
    \cong
    \chi_W(\overline H(U)),
\]
where
\[
    \overline H:[X,\underline{\mathbb V}]
    \to
    [K,\underline{\mathbb V}]
\]
is given by
\[
    \overline H(U)(k)=\chi_{Hk}(U).
\]

The assignment \(I\mapsto\chi_I\) is functorial on weights: a morphism
\(I\to I'\) in the presheaf category induces, by functoriality of the coend, a
morphism
\[
    \chi_I(U)\to\chi_{I'}(U)
\]
natural in \(U\).  Since
\[
    H:K\to\operatorname{Idl}_{\mathbb V}(X)
\]
is a \(\mathbb V\)-functor, the assignment
\[
    k\longmapsto \chi_{Hk}(U)
\]
is an upper predicate on \(K\).  Hence \(\overline H(U)\) is well-defined.

Each \(Hk\) is finite-flat, so each
\[
    \chi_{Hk}:[X,\underline{\mathbb V}]\to\underline{\mathbb V}
\]
preserves the finite limits of the fuzzy-frame doctrine.  Since those finite
limits in \([K,\underline{\mathbb V}]\) are computed pointwise, \(\overline H\)
preserves them.  Since \(W\) is finite-flat, \(\chi_W\) preserves them.  Hence
\[
    \chi_J\cong\chi_W\overline H
\]
preserves the finite limits of the fuzzy-frame doctrine.  Therefore \(J\) is
finite-flat.
\end{proof}

\begin{corollary}[Flattening is finite-flat]
\label{cor:flattening-finite-flat-fdom-rewrite}
If
\[
    \mathcal I:
    \operatorname{Idl}_{\mathbb V}(X)^{\op}
    \to
    \underline{\mathbb V}
\]
is finite-flat, then the flattened weight
\[
    \mu_X(\mathcal I)(x)
    =
    \int^{I\in\operatorname{Idl}_{\mathbb V}(X)}
        \mathcal I(I)\tensor I(x)
\]
is finite-flat.
\end{corollary}

\begin{proof}
Apply Proposition~\ref{prop:finite-flat-saturation-fdom-rewrite} to the identity
diagram
\[
    1_{\operatorname{Idl}_{\mathbb V}(X)}:
    \operatorname{Idl}_{\mathbb V}(X)
    \to
    \operatorname{Idl}_{\mathbb V}(X).
\]
\end{proof}

\section{The KZ property of the ideal monad}
\label{sec:kz-property-ideal-monad-fdom-rewrite}

The preceding chapter constructed the fuzzy ideal monad.  We now prove the
property of that monad needed to identify its strict algebras with finite-flat
colimit completions.

\begin{proposition}[The KZ comparison for the ideal monad]
\label{prop:ideal-monad-kz-fdom-rewrite}
For every small \(\mathbb V\)-category \(X\), there is an enriched natural
transformation
\[
    T_{\mathbb V}\eta_X
    \Longrightarrow
    \eta_{T_{\mathbb V}X}.
\]
Equivalently, for every finite-flat weight
\[
    I:X^{\op}\to\underline{\mathbb V},
\]
there is a canonical morphism in
\[
    [\operatorname{Idl}_{\mathbb V}(X)^{\op},\underline{\mathbb V}]
\]
from the weight \((T_{\mathbb V}\eta_X)(I)\) to the representable weight
\[
    \operatorname{Idl}_{\mathbb V}(X)(-,I).
\]
These transformations satisfy the lax-idempotent, or KZ, coherence equations for
the ideal monad.
\end{proposition}

\begin{proof}
Let
\[
    I\in\operatorname{Idl}_{\mathbb V}(X)
\]
and
\[
    J\in\operatorname{Idl}_{\mathbb V}(X).
\]
The value of \((T_{\mathbb V}\eta_X)(I)\) at \(J\) is
\[
    (T_{\mathbb V}\eta_X)(I)(J)
    =
    \int^{x\in X}
        I(x)\tensor
        \operatorname{Idl}_{\mathbb V}(X)(J,X(-,x)).
\]
By enriched Yoneda,
\[
    I(x)
    \cong
    [X^{\op},\underline{\mathbb V}](X(-,x),I).
\]
Composition in the presheaf category gives morphisms
\[
    \operatorname{Idl}_{\mathbb V}(X)(J,X(-,x))
    \tensor
    [X^{\op},\underline{\mathbb V}](X(-,x),I)
    \longrightarrow
    \operatorname{Idl}_{\mathbb V}(X)(J,I).
\]
Tensoring with the Yoneda identification and integrating over \(x\) gives
\[
    (T_{\mathbb V}\eta_X)(I)(J)
    \longrightarrow
    \operatorname{Idl}_{\mathbb V}(X)(J,I).
\]
This construction is natural in \(J\), and therefore defines a morphism of
weights
\[
    (T_{\mathbb V}\eta_X)(I)
    \longrightarrow
    \operatorname{Idl}_{\mathbb V}(X)(-,I).
\]
It is also enriched-natural in \(I\), because it is induced by composition in the
enriched presheaf category.

This is precisely the standard KZ comparison for a free
\(\Phi_{\mathrm{ff}}\)-cocompletion monad.  The functor \(T_{\mathbb V}\eta_X\)
sends a finite-flat weight to the finite-flat colimit of its principal
approximants, while \(\eta_{T_{\mathbb V}X}\) sends it to its representable
presheaf.  The comparison is induced by the universal cocone from principal
approximants into the weight.

It remains only to identify the coherence.  The comparison just constructed is
induced entirely by composition in
\[
    [X^{\op},\underline{\mathbb V}],
\]
while the multiplication
\[
    \mu_X:
    \operatorname{Idl}_{\mathbb V}(\operatorname{Idl}_{\mathbb V}(X))
    \to
    \operatorname{Idl}_{\mathbb V}(X)
\]
is flattening.  Hence the unit and multiplication compatibilities reduce to
associativity and unitality of composition in the presheaf category, together
with Fubini for the coends defining flattening.

Explicitly, after evaluating at
\[
    J\in\operatorname{Idl}_{\mathbb V}(X),
\]
the two multiplication-coherence composites are both the coend-composite
\[
    \int^{I,x}
        \mathcal I(I)\tensor I(x)\tensor
        \operatorname{Idl}_{\mathbb V}(X)(J,X(-,x))
    \longrightarrow
    \operatorname{Idl}_{\mathbb V}(X)(J,\mu_X\mathcal I),
\]
obtained by first composing
\[
    \operatorname{Idl}_{\mathbb V}(X)(J,X(-,x))
    \tensor
    [X^{\op},\underline{\mathbb V}](X(-,x),I)
    \longrightarrow
    \operatorname{Idl}_{\mathbb V}(X)(J,I),
\]
and then composing with the structure map expressing the flattening
\[
    \mathcal I(I)\tensor I
    \longrightarrow
    \mu_X\mathcal I.
\]
The unit-coherence equations are the corresponding special cases with
representable weights, where the assertion is the enriched Yoneda identity.
Therefore the displayed enriched natural transformation is the KZ comparison for
the ideal monad.
\end{proof}

\begin{theorem}[The algebra structure is left adjoint to Yoneda]
\label{thm:algebra-map-left-adjoint-y-fdom-rewrite}
Let
\[
    \alpha:\operatorname{Idl}_{\mathbb V}(X)\to X
\]
be a strict algebra structure for the ideal monad.  Then
\[
    \alpha\dashv\eta_X
\]
with identity counit
\[
    \alpha\eta_X=1_X.
\]
Consequently, enriched-naturally in
\(I\in\operatorname{Idl}_{\mathbb V}(X)\) and \(x\in X\), there is an
isomorphism in \(\mathbb V\)
\[
    X(\alpha I,x)
    \cong
    \operatorname{Idl}_{\mathbb V}(X)(I,X(-,x)).
\]
\end{theorem}

\begin{proof}
The algebra unit law gives
\[
    \alpha\eta_X=1_X,
\]
so the counit of the desired adjunction is the identity.

It remains to construct the unit
\[
    1_{\operatorname{Idl}_{\mathbb V}(X)}
    \longrightarrow
    \eta_X\alpha.
\]
Using Proposition~\ref{prop:ideal-monad-kz-fdom-rewrite}, we have an enriched
natural transformation
\[
    T_{\mathbb V}\eta_X
    \Longrightarrow
    \eta_{T_{\mathbb V}X}.
\]
Applying \(T_{\mathbb V}\alpha\) gives
\[
    T_{\mathbb V}\alpha\;T_{\mathbb V}\eta_X
    \longrightarrow
    T_{\mathbb V}\alpha\;\eta_{T_{\mathbb V}X}.
\]
The left-hand side is
\[
    T_{\mathbb V}(\alpha\eta_X)
    =
    T_{\mathbb V}(1_X)
    =
    1_{\operatorname{Idl}_{\mathbb V}(X)}.
\]
By naturality of \(\eta\), the right-hand side is
\[
    \eta_X\alpha.
\]
Thus we obtain
\[
    1_{\operatorname{Idl}_{\mathbb V}(X)}
    \longrightarrow
    \eta_X\alpha.
\]

By the KZ coherence established in
Proposition~\ref{prop:ideal-monad-kz-fdom-rewrite}, the above unit and the
identity counit satisfy the triangle identities.  Therefore
\[
    \alpha\dashv\eta_X.
\]
The displayed hom-isomorphism is the hom-form of this enriched adjunction.
\end{proof}

\section{Fuzzy domains as finite-flat-cocomplete categories}
\label{sec:fuzzy-domains-as-ff-cocomplete-rewrite}

\begin{definition}[Fuzzy domain]
\label{def:fuzzy-domain-category-rewrite}
A fuzzy domain is a strict algebra
\[
    \alpha_X:
    \operatorname{Idl}_{\mathbb V}(X)\to X
\]
for the ideal monad.  Thus
\[
    \alpha_X\eta_X=1_X,
    \qquad
    \alpha_X\mu_X
    =
    \alpha_XT_{\mathbb V}\alpha_X.
\]

A fuzzy-domain morphism
\[
    f:(X,\alpha_X)\to(Y,\alpha_Y)
\]
is a \(\mathbb V\)-functor
\[
    f:X\to Y
\]
such that
\[
    f\alpha_X
    =
    \alpha_YT_{\mathbb V}f.
\]
The ordinary category of fuzzy domains and fuzzy-domain morphisms is denoted
\[
    \mathbf{FDom}_{\mathbb V}.
\]
\end{definition}

\begin{definition}[Strict coherent finite-flat colimits]
\label{def:strict-coherent-ff-colimits-rewrite}
A strict coherent choice of finite-flat weighted colimits in a small
\(\mathbb V\)-category \(X\) consists of the following data.

For every small \(\mathbb V\)-category \(K\), every finite-flat weight
\[
    W:K^{\op}\to\underline{\mathbb V},
\]
and every diagram
\[
    D:K\to X,
\]
there is a chosen object
\[
    W*D\in X
\]
with enriched-natural representing isomorphisms
\[
    X(W*D,x)
    \cong
    [K^{\op},\underline{\mathbb V}](W,X(D-,x)).
\]
These choices are strict on representables:
\[
    K(-,k)*D=Dk,
\]
and their iterated-colimit coherence is the flattening coherence
\[
    \alpha_X\mu_X=\alpha_XT_{\mathbb V}(\alpha_X),
\]
where
\[
    \alpha_X(I)=I*1_X.
\]
\end{definition}

\begin{theorem}[Fuzzy domains and finite-flat colimits]
\label{thm:fuzzy-domains-ff-colimits-rewrite}
Let \(X\) be a small \(\mathbb V\)-category.

To give \(X\) the structure of a fuzzy domain is equivalently to give a strict
coherent choice of all finite-flat weighted colimits in \(X\).  Under this
equivalence, the algebra map
\[
    \alpha_X:\operatorname{Idl}_{\mathbb V}(X)\to X
\]
sends a finite-flat weight \(I:X^{\op}\to\underline{\mathbb V}\) to the
\(I\)-weighted colimit of the identity diagram:
\[
    \alpha_X(I)=I*1_X.
\]

More generally, if
\[
    D:K\to X
\]
is a diagram and
\[
    W:K^{\op}\to\underline{\mathbb V}
\]
is finite-flat, then
\[
    W*D
    =
    \alpha_X(D_!W),
\]
where
\[
    (D_!W)(x)
    =
    \int^{k\in K}W(k)\tensor X(x,Dk).
\]

A \(\mathbb V\)-functor between fuzzy domains is a fuzzy-domain morphism if and
only if it preserves finite-flat weighted colimits.
\end{theorem}

\begin{proof}
First suppose that
\[
    \alpha_X:\operatorname{Idl}_{\mathbb V}(X)\to X
\]
is a strict algebra.  By
Theorem~\ref{thm:algebra-map-left-adjoint-y-fdom-rewrite},
\[
    \alpha_X\dashv\eta_X.
\]
Hence, for every finite-flat weight \(I\) and every \(x\in X\),
\[
\begin{aligned}
    X(\alpha_XI,x)
    &\cong
    \operatorname{Idl}_{\mathbb V}(X)(I,\eta_Xx)
    \\
    &=
    [X^{\op},\underline{\mathbb V}](I,X(-,x)).
\end{aligned}
\]
This is precisely the universal property of the \(I\)-weighted colimit of the
identity diagram.  Thus
\[
    \alpha_X(I)=I*1_X.
\]

Now let \(D:K\to X\), and let \(W:K^{\op}\to\underline{\mathbb V}\) be
finite-flat.  By Lemma~\ref{lem:finite-flat-direct-images-fdom-rewrite},
\(D_!W\) is finite-flat.  For every \(x\in X\),
\[
\begin{aligned}
    X(\alpha_X(D_!W),x)
    &\cong
    [X^{\op},\underline{\mathbb V}](D_!W,X(-,x))
    \\
    &\cong
    [K^{\op},\underline{\mathbb V}](W,X(D-,x)).
\end{aligned}
\]
The second isomorphism is the weighted left-Kan-extension adjunction defining
\(D_!W\).  This is the universal property of the weighted colimit \(W*D\).
Therefore
\[
    W*D=\alpha_X(D_!W)
\]
under the strictification convention.

Conversely, suppose \(X\) is equipped with a strict coherent choice of all
finite-flat weighted colimits.  Define
\[
    \alpha_X(I)=I*1_X.
\]
The representing isomorphisms
\[
    X(\alpha_XI,x)
    \cong
    [X^{\op},\underline{\mathbb V}](I,X(-,x))
\]
are enriched-natural in \(I\) and \(x\).  Hence the assignment
\(I\mapsto\alpha_XI\) extends uniquely to a \(\mathbb V\)-functor
\[
    \alpha_X:\operatorname{Idl}_{\mathbb V}(X)\to X
\]
left adjoint to the Yoneda embedding.

For a representable weight,
\[
    X(-,x)*1_X=x,
\]
so
\[
    \alpha_X\eta_X=1_X.
\]
For a finite-flat weight
\[
    \mathcal I:
    \operatorname{Idl}_{\mathbb V}(X)^{\op}
    \to
    \underline{\mathbb V},
\]
the left-hand side of the multiplication law is
\[
    \alpha_X\mu_X(\mathcal I)
    =
    \mu_X(\mathcal I)*1_X.
\]
The right-hand side is
\[
    \alpha_XT_{\mathbb V}\alpha_X(\mathcal I)
    =
    (T_{\mathbb V}\alpha_X)(\mathcal I)*1_X.
\]
The weight \(T_{\mathbb V}\alpha_X(\mathcal I)\) on \(X\) is the lower direct
image of \(\mathcal I\) along
\[
    \alpha_X:\operatorname{Idl}_{\mathbb V}(X)\to X.
\]
Therefore its colimit against \(1_X\) is the same weighted colimit as
\(\mathcal I\) against the diagram \(\alpha_X\):
\[
    (T_{\mathbb V}\alpha_X)(\mathcal I)*1_X
    \cong
    \mathcal I*\alpha_X.
\]
But
\[
    \mu_X(\mathcal I)(x)
    =
    \int^{I}\mathcal I(I)\tensor I(x)
\]
is the flattened weight computing the same iterated finite-flat colimit.  By the
strict associativity coherence of the chosen finite-flat colimits,
\[
    \mu_X(\mathcal I)*1_X
    =
    \mathcal I*(I\mapsto I*1_X)
    =
    \mathcal I*\alpha_X.
\]
Thus
\[
    \alpha_X\mu_X
    =
    \alpha_XT_{\mathbb V}\alpha_X,
\]
so \(\alpha_X\) is a strict algebra.

Finally, let
\[
    f:(X,\alpha_X)\to(Y,\alpha_Y)
\]
be a \(\mathbb V\)-functor.  If \(f\) is a fuzzy-domain morphism, then for every
finite-flat \(I:X^{\op}\to\underline{\mathbb V}\),
\[
    f(I*1_X)
    =
    f\alpha_X(I)
    =
    \alpha_YT_{\mathbb V}f(I)
    =
    T_{\mathbb V}f(I)*1_Y.
\]
Thus \(f\) preserves finite-flat colimits of identity diagrams, and hence, by
the direct-image formula, all finite-flat weighted colimits.

Conversely, if \(f\) preserves finite-flat weighted colimits, then for every
finite-flat \(I\),
\[
    f\alpha_X(I)
    =
    f(I*1_X)
    =
    T_{\mathbb V}f(I)*1_Y
    =
    \alpha_YT_{\mathbb V}f(I).
\]
Therefore
\[
    f\alpha_X=\alpha_YT_{\mathbb V}f.
\]
Thus \(f\) is a fuzzy-domain morphism.
\end{proof}

\section{Free fuzzy domains and limits}
\label{sec:free-fuzzy-domains-limits-rewrite}

\begin{theorem}[Free fuzzy domains]
\label{thm:free-fuzzy-domains-rewrite}
The forgetful functor
\[
    U:\mathbf{FDom}_{\mathbb V}\to\mathbb V\text{-}\mathbf{Cat}
\]
has a left adjoint
\[
    F:\mathbb V\text{-}\mathbf{Cat}\to\mathbf{FDom}_{\mathbb V},
    \qquad
    F(X)=\operatorname{Idl}_{\mathbb V}(X),
\]
where \(F(X)\) has algebra structure
\[
    \mu_X:
    \operatorname{Idl}_{\mathbb V}(\operatorname{Idl}_{\mathbb V}(X))
    \to
    \operatorname{Idl}_{\mathbb V}(X).
\]
Thus, for every small \(\mathbb V\)-category \(X\) and every fuzzy domain \(A\),
there is a natural bijection
\[
    \mathbf{FDom}_{\mathbb V}
    \bigl(
        \operatorname{Idl}_{\mathbb V}(X),
        A
    \bigr)
    \cong
    \mathbb V\text{-}\mathbf{Cat}(X,UA).
\]
\end{theorem}

\begin{proof}
This is the free-algebra adjunction of the monad \(T_{\mathbb V}\).  Explicitly,
if
\[
    f:X\to UA
\]
is a \(\mathbb V\)-functor and
\[
    \alpha_A:\operatorname{Idl}_{\mathbb V}(A)\to A
\]
is the algebra map of \(A\), define
\[
    \overline f:
    \operatorname{Idl}_{\mathbb V}(X)\to A
\]
by
\[
    \overline f(I)
    =
    \alpha_A(T_{\mathbb V}f(I)).
\]
Equivalently,
\[
    \overline f(I)=I*f.
\]
By Theorem~\ref{thm:fuzzy-domains-ff-colimits-rewrite},
\(\overline f\) preserves finite-flat weighted colimits, so it is a
fuzzy-domain morphism.  Moreover,
\[
    \overline f(\eta_Xx)
    =
    \alpha_A(\eta_A(fx))
    =
    fx.
\]
Thus
\[
    \overline f\eta_X=f.
\]

Conversely, if
\[
    g:\operatorname{Idl}_{\mathbb V}(X)\to A
\]
is a fuzzy-domain morphism, then \(g\) preserves finite-flat weighted colimits.
For every finite-flat weight \(I:X^{\op}\to\underline{\mathbb V}\), the object
\(I\in\operatorname{Idl}_{\mathbb V}(X)\) is the \(I\)-weighted colimit of the
Yoneda diagram:
\[
    I \cong I*\eta_X.
\]
Therefore
\[
    g(I)
    \cong
    g(I*\eta_X)
    \cong
    I*(g\eta_X).
\]
Hence
\[
    g=\overline{g\eta_X}
\]
under the strictification convention.  The two constructions are inverse and
natural.
\end{proof}

\begin{theorem}[Monadicity]
\label{thm:fdom-monadicity-rewrite}
The forgetful functor
\[
    U:\mathbf{FDom}_{\mathbb V}\to\mathbb V\text{-}\mathbf{Cat}
\]
is monadic, and the induced monad is the fuzzy ideal monad
\[
    T_{\mathbb V}.
\]
Thus
\[
    \mathbf{FDom}_{\mathbb V}
    =
    (\mathbb V\text{-}\mathbf{Cat})^{T_{\mathbb V}}.
\]
\end{theorem}

\begin{proof}
By definition, \(\mathbf{FDom}_{\mathbb V}\) is the Eilenberg--Moore category of
strict algebras for \(T_{\mathbb V}\).  The adjunction of
Theorem~\ref{thm:free-fuzzy-domains-rewrite} has underlying monad
\[
    UF=T_{\mathbb V},
\]
with Yoneda unit and flattening multiplication.  Therefore \(U\) is monadic.
\end{proof}

\begin{theorem}[Small conical limits]
\label{thm:fdom-limits-rewrite}
The category
\[
    \mathbf{FDom}_{\mathbb V}
\]
has all small conical limits whose underlying limits exist in
\[
    \mathbb V\text{-}\mathbf{Cat}.
\]
These limits are created by the forgetful functor
\[
    U:\mathbf{FDom}_{\mathbb V}\to\mathbb V\text{-}\mathbf{Cat}.
\]
\end{theorem}

\begin{proof}
Eilenberg--Moore categories have all limits that exist in the base category, and
the forgetful functor creates them.

Explicitly, let
\[
    D:J\to\mathbf{FDom}_{\mathbb V}
\]
be a small diagram, and suppose that
\[
    L=\lim_{j\in J}UD_j
\]
exists in \(\mathbb V\text{-}\mathbf{Cat}\), with projections
\[
    p_j:L\to UD_j.
\]
There is a unique algebra structure
\[
    \alpha_L:T_{\mathbb V}L\to L
\]
such that
\[
    p_j\alpha_L
    =
    \alpha_jT_{\mathbb V}(p_j)
\]
for every \(j\).  With this algebra structure, \(L\) satisfies the universal
property of the limit in \(\mathbf{FDom}_{\mathbb V}\).
\end{proof}

\section{External products and Fubini}
\label{sec:external-products-fubini-rewrite}

\begin{definition}[External product of weights]
\label{def:external-product-weights-rewrite}
Let
\[
    W:K^{\op}\to\underline{\mathbb V},
    \qquad
    V:J^{\op}\to\underline{\mathbb V}
\]
be weights.  Their external product is the weight
\[
    W\boxtimes V:(K\tensor J)^{\op}\to\underline{\mathbb V}
\]
defined by
\[
    (W\boxtimes V)(k,j)=W(k)\tensor V(j).
\]
\end{definition}

\begin{proposition}[External products of finite-flat weights]
\label{prop:external-products-finite-flat-rewrite}
If \(W\) and \(V\) are finite-flat, then
\[
    W\boxtimes V
\]
is finite-flat.
\end{proposition}

\begin{proof}
Let
\[
    E:K\tensor J\to\underline{\mathbb V}
\]
be an upper predicate.  Then
\[
\begin{aligned}
    \chi_{W\boxtimes V}(E)
    &=
    \int^{(k,j)\in K\tensor J}
        W(k)\tensor V(j)\tensor E(k,j)
    \\
    &\cong
    \int^{k\in K}
        W(k)\tensor
        \left(
            \int^{j\in J}
                V(j)\tensor E(k,j)
        \right).
\end{aligned}
\]
For each \(k\), the inner expression is
\[
    \chi_V(E(k,-)).
\]
Since \(V\) is finite-flat, this construction preserves the finite limits of the
fuzzy-frame doctrine in the \(J\)-variable.

The assignment
\[
    k\longmapsto \chi_V(E(k,-))
\]
is an upper predicate on \(K\), by functoriality of the character
\(V\mapsto\chi_V\).  Since \(W\) is finite-flat, applying \(\chi_W\) preserves
the finite limits of the fuzzy-frame doctrine.  Hence \(\chi_{W\boxtimes V}\)
preserves those finite limits, so \(W\boxtimes V\) is finite-flat.
\end{proof}

\begin{proposition}[Finite-flat Fubini]
\label{prop:finite-flat-fubini-rewrite}
Let
\[
    W:K^{\op}\to\underline{\mathbb V},
    \qquad
    V:J^{\op}\to\underline{\mathbb V}
\]
be finite-flat.  For every diagram
\[
    E:K\tensor J\to C
\]
in a fuzzy domain \(C\), there are canonical isomorphisms
\[
    W*\bigl(k\mapsto V*(j\mapsto E(k,j))\bigr)
    \cong
    (W\boxtimes V)*E
    \cong
    V*\bigl(j\mapsto W*(k\mapsto E(k,j))\bigr).
\]
\end{proposition}

\begin{proof}
These are the Fubini isomorphisms for enriched weighted colimits.  By
Proposition~\ref{prop:external-products-finite-flat-rewrite}, the external
product \(W\boxtimes V\) is finite-flat.  Since \(C\) is a fuzzy domain, all
three finite-flat weighted colimits appearing above exist.
\end{proof}

\section{Internal homs}
\label{sec:internal-homs-fdom-rewrite}

Let \(B\) and \(C\) be fuzzy domains.  The enriched functor category
\[
    [B,C]
\]
contains a full \(\mathbb V\)-subcategory spanned by those
\(\mathbb V\)-functors \(B\to C\) that preserve finite-flat weighted colimits.

\begin{definition}[Internal hom fuzzy domain]
\label{def:internal-hom-fdom-rewrite}
For fuzzy domains \(B\) and \(C\), define
\[
    [B,C]_{\mathrm{ff}}
\]
to be the full \(\mathbb V\)-subcategory of \([B,C]\) spanned by the
finite-flat-colimit-preserving \(\mathbb V\)-functors
\[
    B\to C.
\]
\end{definition}

\begin{proposition}[Internal homs are fuzzy domains]
\label{prop:internal-homs-fdom-rewrite}
For fuzzy domains \(B\) and \(C\), the \(\mathbb V\)-category
\[
    [B,C]_{\mathrm{ff}}
\]
is a fuzzy domain.  Its finite-flat weighted colimits are computed pointwise in
\(C\).
\end{proposition}

\begin{proof}
Let
\[
    W:K^{\op}\to\underline{\mathbb V}
\]
be finite-flat, and let
\[
    H:K\to[B,C]_{\mathrm{ff}}
\]
be a diagram.  Since \(C\) is a fuzzy domain, the pointwise formula
\[
    (W*H)(b)
    =
    W*(k\mapsto Hk(b))
\]
defines the \(W\)-weighted colimit of \(H\) in the ambient functor category
\([B,C]\).

It remains to show that \(W*H\) preserves finite-flat weighted colimits in \(B\).
Let
\[
    V:J^{\op}\to\underline{\mathbb V}
\]
be finite-flat, and let
\[
    D:J\to B
\]
be a diagram.  Then
\[
\begin{aligned}
    (W*H)(V*D)
    &\cong
    W*\bigl(k\mapsto Hk(V*D)\bigr)
    \\
    &\cong
    W*\bigl(k\mapsto V*(j\mapsto Hk(Dj))\bigr)
    \\
    &\cong
    V*\bigl(j\mapsto W*(k\mapsto Hk(Dj))\bigr)
    \\
    &\cong
    V*((W*H)D).
\end{aligned}
\]
The second isomorphism uses that each
\[
    Hk:B\to C
\]
preserves finite-flat weighted colimits.  The third isomorphism is
Proposition~\ref{prop:finite-flat-fubini-rewrite}.  Hence \(W*H\) preserves
finite-flat weighted colimits in \(B\).

Therefore \([B,C]_{\mathrm{ff}}\) is closed under finite-flat weighted colimits
inside \([B,C]\).  Since \([B,C]_{\mathrm{ff}}\) is full in \([B,C]\), the
ambient weighted-colimit universal property restricts to
\([B,C]_{\mathrm{ff}}\).  By
Theorem~\ref{thm:fuzzy-domains-ff-colimits-rewrite},
\([B,C]_{\mathrm{ff}}\) is a fuzzy domain.

The algebra structure on \([B,C]_{\mathrm{ff}}\) sends a finite-flat weight
\(W\) on a diagram \(H\) to the pointwise colimit
\[
    b\longmapsto W*(k\mapsto Hk(b)).
\]
\end{proof}

\section{Bimorphisms and the Kelly tensor product}
\label{sec:bimorphisms-kelly-tensor-rewrite}

\begin{definition}[Finite-flat bimorphism]
\label{def:finite-flat-bimorphism-rewrite}
Let \(A,B,C\) be fuzzy domains.  A finite-flat bimorphism
\[
    M:A\tensor B\to C
\]
is a \(\mathbb V\)-functor such that:
\begin{enumerate}
    \item for each \(a\in A\), the partial functor
    \[
        M(a,-):B\to C
    \]
    preserves finite-flat weighted colimits;
    \item for each \(b\in B\), the partial functor
    \[
        M(-,b):A\to C
    \]
    preserves finite-flat weighted colimits.
\end{enumerate}
We write
\[
    \operatorname{Bimor}_{\mathrm{ff}}(A,B;C)
\]
for the set of finite-flat bimorphisms from \(A,B\) to \(C\).
\end{definition}

\begin{proposition}[Currying finite-flat bimorphisms]
\label{prop:currying-bimorphisms-rewrite}
For fuzzy domains \(A,B,C\), there is a natural bijection
\[
    \operatorname{Bimor}_{\mathrm{ff}}(A,B;C)
    \cong
    \mathbf{FDom}_{\mathbb V}
    \bigl(A,[B,C]_{\mathrm{ff}}\bigr).
\]
\end{proposition}

\begin{proof}
Let
\[
    M:A\tensor B\to C
\]
be a finite-flat bimorphism.  By enriched currying, \(M\) corresponds to a
\(\mathbb V\)-functor
\[
    \widehat M:A\to[B,C],
    \qquad
    \widehat M(a)(b)=M(a,b).
\]
Since \(M\) preserves finite-flat weighted colimits in the \(B\)-variable, each
\[
    \widehat M(a):B\to C
\]
lies in \([B,C]_{\mathrm{ff}}\).  Thus \(\widehat M\) factors through
\[
    [B,C]_{\mathrm{ff}}\subseteq[B,C].
\]
Since \(M\) preserves finite-flat weighted colimits in the \(A\)-variable and
finite-flat colimits in \([B,C]_{\mathrm{ff}}\) are computed pointwise,
\[
    \widehat M:A\to[B,C]_{\mathrm{ff}}
\]
preserves finite-flat weighted colimits.  Hence \(\widehat M\) is a
fuzzy-domain morphism.

Conversely, let
\[
    H:A\to[B,C]_{\mathrm{ff}}
\]
be a fuzzy-domain morphism.  Uncurrying gives
\[
    \widetilde H:A\tensor B\to C,
    \qquad
    \widetilde H(a,b)=H(a)(b).
\]
For each \(a\), the partial functor
\[
    \widetilde H(a,-)=H(a)
\]
preserves finite-flat weighted colimits because \(H(a)\) is an object of
\([B,C]_{\mathrm{ff}}\).  For each \(b\), the partial functor
\[
    \widetilde H(-,b):A\to C
\]
preserves finite-flat weighted colimits because \(H\) does and colimits in
\([B,C]_{\mathrm{ff}}\) are computed pointwise.  Hence \(\widetilde H\) is a
finite-flat bimorphism.

The two constructions are inverse by the enriched tensor--hom adjunction.
\end{proof}

The preceding saturation and Fubini results supply the hypotheses needed for
Kelly's tensor-product theorem for categories with a saturated commutative class
of specified weighted colimits.  We apply that theorem to the doctrine
\[
    \Phi_{\mathrm{ff}}
\]
of finite-flat weights.  All constructions are interpreted under the universe
conventions fixed in the earlier chapters.

\begin{theorem}[Kelly tensor product for fuzzy domains]
\label{thm:kelly-tensor-product-fdom-rewrite}
For fuzzy domains \(A\) and \(B\), there is a fuzzy domain
\[
    A\boxtimes_{\mathrm{ff}}B
\]
and a finite-flat bimorphism
\[
    \kappa_{A,B}:A\tensor B\to A\boxtimes_{\mathrm{ff}}B
\]
such that, for every fuzzy domain \(C\), composition with \(\kappa_{A,B}\)
induces a natural bijection
\[
    \mathbf{FDom}_{\mathbb V}
    (A\boxtimes_{\mathrm{ff}}B,C)
    \cong
    \operatorname{Bimor}_{\mathrm{ff}}(A,B;C).
\]
\end{theorem}

\begin{proof}
By Lemma~\ref{lem:representables-finite-flat-fdom-rewrite},
Lemma~\ref{lem:finite-flat-direct-images-fdom-rewrite}, and
Proposition~\ref{prop:finite-flat-saturation-fdom-rewrite}, the assignment
\[
    X\longmapsto \Phi_{\mathrm{ff}}(X)
\]
contains the representables, is stable under lower direct image along
\(\mathbb V\)-functors, and is closed under \(\Phi_{\mathrm{ff}}\)-weighted
colimits in presheaf categories.  Hence it is saturated in Kelly's sense.

Moreover, Proposition~\ref{prop:external-products-finite-flat-rewrite} and
Proposition~\ref{prop:finite-flat-fubini-rewrite} show that
\(\Phi_{\mathrm{ff}}\) is a commutative doctrine: finite-flat weighted colimits
commute with finite-flat weighted colimits.

Therefore Kelly's tensor-product theorem for categories equipped with a
saturated commutative class of specified weighted colimits applies to the
finite-flat doctrine.  The representing object for separately
finite-flat-colimit-preserving functors
\[
    A\tensor B\to C
\]
is denoted
\[
    A\boxtimes_{\mathrm{ff}}B.
\]
Its universal map
\[
    \kappa_{A,B}:A\tensor B\to A\boxtimes_{\mathrm{ff}}B
\]
is a finite-flat bimorphism, and the displayed bijection is exactly the
representing universal property.
\end{proof}

\begin{corollary}[Tensor product of free fuzzy domains]
\label{cor:tensor-free-fdom-rewrite}
For small \(\mathbb V\)-categories \(X\) and \(Y\), there is a canonical
equivalence
\[
    \operatorname{Idl}_{\mathbb V}(X)
    \boxtimes_{\mathrm{ff}}
    \operatorname{Idl}_{\mathbb V}(Y)
    \simeq
    \operatorname{Idl}_{\mathbb V}(X\tensor Y).
\]
\end{corollary}

\begin{proof}
For every fuzzy domain \(C\), there are natural bijections
\[
\begin{aligned}
\mathbf{FDom}_{\mathbb V}
\bigl(
    \operatorname{Idl}_{\mathbb V}(X)
    \boxtimes_{\mathrm{ff}}
    \operatorname{Idl}_{\mathbb V}(Y),
    C
\bigr)
&\cong
\operatorname{Bimor}_{\mathrm{ff}}
\bigl(
    \operatorname{Idl}_{\mathbb V}(X),
    \operatorname{Idl}_{\mathbb V}(Y);
    C
\bigr)
\\
&\cong
\mathbb V\text{-}\mathbf{Cat}(X\tensor Y,UC)
\\
&\cong
\mathbf{FDom}_{\mathbb V}
\bigl(
    \operatorname{Idl}_{\mathbb V}(X\tensor Y),
    C
\bigr).
\end{aligned}
\]

We spell out the middle bijection.  A finite-flat bimorphism
\[
    M:\operatorname{Idl}_{\mathbb V}(X)\tensor
      \operatorname{Idl}_{\mathbb V}(Y)\to C
\]
restricts along
\[
    \eta_X\tensor\eta_Y:X\tensor Y\to
    \operatorname{Idl}_{\mathbb V}(X)\tensor
    \operatorname{Idl}_{\mathbb V}(Y)
\]
to a \(\mathbb V\)-functor
\[
    m:X\tensor Y\to UC.
\]
Conversely, given
\[
    m:X\tensor Y\to UC,
\]
define
\[
    \widetilde m:
    \operatorname{Idl}_{\mathbb V}(X)\tensor
    \operatorname{Idl}_{\mathbb V}(Y)
    \to C
\]
by
\[
    \widetilde m(I,J)
    =
    I*\bigl(x\mapsto J*(y\mapsto m(x,y))\bigr).
\]
Since \(I\) and \(J\) are finite-flat, \(I\boxtimes J\) is finite-flat by
Proposition~\ref{prop:external-products-finite-flat-rewrite}.  By finite-flat
Fubini,
\[
    \widetilde m(I,J)
    \cong
    (I\boxtimes J)*m.
\]
The functor \(\widetilde m\) preserves finite-flat weighted colimits separately
in \(I\) and \(J\), and it is the unique finite-flat bimorphism extending
\(m\).  This gives the middle bijection.

By representability, the displayed equivalence follows.
\end{proof}

\begin{corollary}[Tensor unit]
\label{cor:tensor-unit-fdom-rewrite}
Let
\[
    \mathbb 1
\]
be the terminal small \(\mathbb V\)-category.  The tensor unit for fuzzy domains
is
\[
    \operatorname{Idl}_{\mathbb V}(\mathbb 1).
\]
\end{corollary}

\begin{proof}
For every fuzzy domain \(C\),
\[
\begin{aligned}
    \mathbf{FDom}_{\mathbb V}
    \bigl(
        \operatorname{Idl}_{\mathbb V}(\mathbb 1)
        \boxtimes_{\mathrm{ff}}A,
        C
    \bigr)
    &\cong
    \operatorname{Bimor}_{\mathrm{ff}}
    \bigl(
        \operatorname{Idl}_{\mathbb V}(\mathbb 1),
        A;
        C
    \bigr)
    \\
    &\cong
    \mathbf{FDom}_{\mathbb V}(A,C).
\end{aligned}
\]
The second bijection uses that
\(\operatorname{Idl}_{\mathbb V}(\mathbb 1)\) is freely generated as a fuzzy
domain by the unique object of \(\mathbb 1\).  Hence
\[
    \operatorname{Idl}_{\mathbb V}(\mathbb 1)
    \boxtimes_{\mathrm{ff}}A
    \simeq
    A.
\]
The right unit law is analogous.
\end{proof}

\section{Symmetric monoidal closure}
\label{sec:smc-fdom-rewrite}

\begin{theorem}[Symmetric monoidal closure of fuzzy domains]
\label{thm:fdom-smc-rewrite}
The category
\[
    \mathbf{FDom}_{\mathbb V}
\]
is symmetric monoidal closed.  Its tensor product is
\[
    A\boxtimes_{\mathrm{ff}}B,
\]
its tensor unit is
\[
    \operatorname{Idl}_{\mathbb V}(\mathbb 1),
\]
and its internal hom is
\[
    [B,C]_{\mathrm{ff}}.
\]
The closed structure is the natural bijection
\[
    \mathbf{FDom}_{\mathbb V}
    (A\boxtimes_{\mathrm{ff}}B,C)
    \cong
    \mathbf{FDom}_{\mathbb V}
    \bigl(A,[B,C]_{\mathrm{ff}}\bigr).
\]
\end{theorem}

\begin{proof}
By Theorem~\ref{thm:kelly-tensor-product-fdom-rewrite},
\[
    \mathbf{FDom}_{\mathbb V}
    (A\boxtimes_{\mathrm{ff}}B,C)
    \cong
    \operatorname{Bimor}_{\mathrm{ff}}(A,B;C).
\]
By Proposition~\ref{prop:currying-bimorphisms-rewrite},
\[
    \operatorname{Bimor}_{\mathrm{ff}}(A,B;C)
    \cong
    \mathbf{FDom}_{\mathbb V}
    \bigl(A,[B,C]_{\mathrm{ff}}\bigr).
\]
Composing these bijections gives the closed structure.

The symmetry of \(\boxtimes_{\mathrm{ff}}\) follows from the symmetry of
\[
    A\tensor B\cong B\tensor A
\]
and from the symmetric nature of the condition of preserving finite-flat
weighted colimits separately in the two variables.

Associativity follows by representability.  For every fuzzy domain \(D\), both
\[
    (A\boxtimes_{\mathrm{ff}}B)\boxtimes_{\mathrm{ff}}C
\]
and
\[
    A\boxtimes_{\mathrm{ff}}(B\boxtimes_{\mathrm{ff}}C)
\]
represent trilinear \(\mathbb V\)-functors
\[
    A\tensor B\tensor C\to D
\]
that preserve finite-flat weighted colimits separately in each variable.  Hence
they are canonically equivalent.

The tensor unit is
\[
    \operatorname{Idl}_{\mathbb V}(\mathbb 1)
\]
by Corollary~\ref{cor:tensor-unit-fdom-rewrite}.  Therefore
\(\mathbf{FDom}_{\mathbb V}\) is symmetric monoidal closed.
\end{proof}

\begin{corollary}[Evaluation]
\label{cor:evaluation-fdom-rewrite}
For fuzzy domains \(B\) and \(C\), there is an evaluation morphism
\[
    [B,C]_{\mathrm{ff}}
    \boxtimes_{\mathrm{ff}}
    B
    \longrightarrow
    C.
\]
On generators it is given by
\[
    (f,b)\longmapsto f(b).
\]
\end{corollary}

\begin{proof}
The ordinary enriched evaluation
\[
    [B,C]\tensor B\to C
\]
restricts to a finite-flat bimorphism
\[
    [B,C]_{\mathrm{ff}}\tensor B\to C.
\]
It is finite-flat-colimit-preserving in the \(B\)-variable because every object
of \([B,C]_{\mathrm{ff}}\) preserves finite-flat weighted colimits.  It is
finite-flat-colimit-preserving in the \([B,C]_{\mathrm{ff}}\)-variable because
finite-flat colimits there are computed pointwise in \(C\).  By the universal
property of \(\boxtimes_{\mathrm{ff}}\), this bimorphism induces the displayed
fuzzy-domain morphism.
\end{proof}

\section{The posetal specialization}
\label{sec:posetal-specialization-fdom-rewrite}

Let
\[
    \Lambda=(L,\leq,\meet,\join,\tensor,\multimap,\bot_L,\top_L)
\]
be a fuzzy truth-value object.  A fuzzy domain over \(\Lambda\) is a fuzzy
preorder \(X\) equipped with a strict algebra
\[
    \alpha_X:\operatorname{Idl}_{\Lambda}(X)\to X.
\]
The map \(\alpha_X\) sends a fuzzy ideal \(I\) to its fuzzy ideal-join.

If
\[
    D:K\to X
\]
is a \(\Lambda\)-functor and
\[
    W:K^{\op}\to\Lambda
\]
is a fuzzy ideal of \(K\), then
\[
    W*D
    =
    \alpha_X(D_!W),
\]
where
\[
    (D_!W)(x)
    =
    \bigjoin_{k\in K}
        W(k)\tensor X(x,Dk).
\]
A morphism of fuzzy domains
\[
    f:X\to Y
\]
is a \(\Lambda\)-functor satisfying
\[
    f\alpha_X
    =
    \alpha_YT_{\Lambda}f.
\]
Equivalently, for every fuzzy ideal \(I\) of \(X\),
\[
    f(\alpha_X(I))
    =
    \alpha_Y(T_{\Lambda}f(I)).
\]
Here the direct image \(T_\Lambda f(I)\) is given by
\[
    (T_\Lambda f)(I)(y)
    =
    \bigjoin_{x\in X} I(x)\tensor Y(y,fx).
\]

For fuzzy domains \(B\) and \(C\), the internal hom
\[
    [B,C]_{\mathrm{ff}}
\]
has as objects the fuzzy-domain morphisms
\[
    B\to C.
\]
Its hom-values are inherited from the enriched functor preorder:
\[
    [B,C]_{\mathrm{ff}}(f,g)
    =
    \bigmeet_{b\in B}C(fb,gb).
\]
Finite-flat joins in \([B,C]_{\mathrm{ff}}\) are computed pointwise in \(C\).

The tensor product
\[
    A\boxtimes_{\mathrm{ff}}B
\]
represents fuzzy bimorphisms
\[
    A\tensor B\to C
\]
that preserve fuzzy ideal joins separately in each variable.

\section{The crisp separated case}
\label{sec:crisp-separated-case-fdom-rewrite}

For \(\Lambda=2\), finite-flat weights are ordinary ideals.  On separated
preorders, strict algebras for the ideal monad are precisely dcpos, and
fuzzy-domain morphisms are precisely Scott-continuous maps.  Thus
\[
    \mathbf{FDom}_{2}^{\mathrm{sep}}
    \simeq
    \mathbf{DCPO},
\]
where \(\mathbf{DCPO}\) denotes the category of dcpos small in the chosen
universe.

In this case, finite-flat bimorphisms are monotone maps
\[
    A\times B\to C
\]
between dcpos that are Scott-continuous in each variable separately.  The
monotonicity is part of the bimorphism data, since a bimorphism is in particular
a \(2\)-functor
\[
    A\tensor B\to C.
\]
Such maps are Scott-continuous for the product dcpo.

Indeed, let
\[
    D\subseteq A\times B
\]
be nonempty and directed.  Then \(\pi_A D\) and \(\pi_BD\) are nonempty directed
subsets, and
\[
    \bigjoin D
    =
    \left(
        \bigjoin\pi_A D,
        \bigjoin\pi_BD
    \right).
\]
If
\[
    f:A\times B\to C
\]
is separately Scott-continuous, then
\[
\begin{aligned}
    f\left(\bigjoin D\right)
    &=
    f\left(
        \bigjoin\pi_A D,
        \bigjoin\pi_B D
    \right)
    \\
    &=
    \bigjoin_{a\in\pi_A D}
        f\left(a,\bigjoin\pi_BD\right)
    \\
    &=
    \bigjoin_{a\in\pi_A D}
        \bigjoin_{b\in\pi_BD}
            f(a,b).
\end{aligned}
\]
Since \(D\) is directed, every pair
\[
    (a,b)\in\pi_A D\times\pi_BD
\]
is bounded above by some element of \(D\).  By monotonicity,
\[
    \bigjoin_{a\in\pi_A D,\;b\in\pi_BD}f(a,b)
    =
    \bigjoin_{(a,b)\in D}f(a,b).
\]
Hence
\[
    f\left(\bigjoin D\right)
    =
    \bigjoin_{(a,b)\in D}f(a,b).
\]
Thus \(f\) is Scott-continuous on the product dcpo.

Consequently, the product dcpo \(A\times B\) represents finite-flat bimorphisms.
By the universal property of the Kelly tensor product, there is a canonical
isomorphism
\[
    A\boxtimes_{\mathrm{ff}}B
    \cong
    A\times B.
\]
The internal hom is the dcpo of Scott-continuous maps
\[
    B\to C
\]
ordered pointwise, with directed joins computed pointwise.

Thus, under the equivalence
\[
    \mathbf{FDom}_{2}^{\mathrm{sep}}
    \simeq
    \mathbf{DCPO},
\]
the symmetric monoidal closed structure constructed above is identified with the
usual cartesian closed structure of dcpos.

\section{Products versus tensor products}
\label{sec:products-versus-tensors-fdom-rewrite}

The category \(\mathbf{FDom}_{\mathbb V}\) has categorical products whenever the
corresponding products exist in
\[
    \mathbb V\text{-}\mathbf{Cat},
\]
and these products are created by the forgetful functor.  The categorical
product
\[
    A\times B
\]
is characterized by
\[
    \mathbf{FDom}_{\mathbb V}(X,A\times B)
    \cong
    \mathbf{FDom}_{\mathbb V}(X,A)
    \times
    \mathbf{FDom}_{\mathbb V}(X,B).
\]

The Kelly tensor product
\[
    A\boxtimes_{\mathrm{ff}}B
\]
is characterized instead by
\[
    \mathbf{FDom}_{\mathbb V}(A\boxtimes_{\mathrm{ff}}B,C)
    \cong
    \operatorname{Bimor}_{\mathrm{ff}}(A,B;C).
\]
Thus the categorical product represents pairs of morphisms out of a common
domain, whereas the Kelly tensor product represents separately
finite-flat-colimit-preserving bimorphisms.

For general fuzzy truth-value bases these are different universal properties, so
there is no canonical identification between
\[
    A\times B
\]
and
\[
    A\boxtimes_{\mathrm{ff}}B
\]
without a separate comparison theorem.  In the crisp separated case, the
separate theorem is that separately Scott-continuous maps out of the product
dcpo are jointly Scott-continuous.  Hence, for dcpos,
\[
    A\boxtimes_{\mathrm{ff}}B
    \cong
    A\times B.
\]

\section{Summary}
\label{sec:summary-fdom-rewrite}

The category of fuzzy domains is the Eilenberg--Moore category of the fuzzy ideal
monad:
\[
    \mathbf{FDom}_{\mathbb V}
    =
    (\mathbb V\text{-}\mathbf{Cat})^{T_{\mathbb V}}.
\]
A fuzzy domain is equivalently a small \(\mathbb V\)-category equipped with a
strict coherent choice of finite-flat weighted colimits, and a fuzzy-domain
morphism is equivalently a \(\mathbb V\)-functor preserving those colimits.

The forgetful functor
\[
    U:\mathbf{FDom}_{\mathbb V}\to\mathbb V\text{-}\mathbf{Cat}
\]
is monadic.  Its left adjoint sends
\[
    X
    \longmapsto
    \operatorname{Idl}_{\mathbb V}(X).
\]
The category \(\mathbf{FDom}_{\mathbb V}\) has the small conical limits created
by \(U\), whenever the corresponding underlying limits exist in
\(\mathbb V\text{-}\mathbf{Cat}\).

The finite-flat doctrine is saturated and commutative.  It is saturated because
finite-flat weights contain the representables, are stable under lower direct
image, and are closed under finite-flat weighted colimits in presheaf
categories.  It is commutative because external products of finite-flat weights
are finite-flat and finite-flat weighted colimits satisfy Fubini.  Therefore
Kelly's tensor-product theorem applies to the finite-flat doctrine.  The
resulting tensor product
\[
    A\boxtimes_{\mathrm{ff}}B
\]
represents finite-flat bimorphisms
\[
    A\tensor B\to C.
\]
The internal hom is
\[
    [B,C]_{\mathrm{ff}},
\]
the fuzzy domain of finite-flat-colimit-preserving functors
\[
    B\to C.
\]
Thus
\[
    \mathbf{FDom}_{\mathbb V}
    (A\boxtimes_{\mathrm{ff}}B,C)
    \cong
    \mathbf{FDom}_{\mathbb V}
    \bigl(A,[B,C]_{\mathrm{ff}}\bigr),
\]
so \(\mathbf{FDom}_{\mathbb V}\) is symmetric monoidal closed.

In the crisp separated case,
\[
    \mathbf{FDom}_{2}^{\mathrm{sep}}
    \simeq
    \mathbf{DCPO},
\]
where dcpos are taken small in the chosen universe.  The Kelly tensor product is
canonically identified with the cartesian product of dcpos,
\[
    A\boxtimes_{\mathrm{ff}}B
    \cong
    A\times B,
\]
and the internal hom is the dcpo of Scott-continuous maps ordered pointwise.